\documentclass[11pt,a4paper]{article}

% ---------------------------------------------------------------
% Pacchetti (non serve toccarli)
% ---------------------------------------------------------------
\usepackage[T1]{fontenc}
\usepackage[utf8]{inputenc}
\usepackage[italian]{babel}
\usepackage[a4paper,margin=2.5cm]{geometry}
\usepackage{amsmath}
\usepackage{booktabs}
\usepackage{graphicx}
\usepackage{enumitem}
\usepackage{url}


\setlist{itemsep=2pt, topsep=4pt}

% Spazio vuoto da riempire: nel testo compare come una riga.
% Per trovarli tutti, cerca \vuoto in questo file.
\newcommand{\vuoto}[1][2.5cm]{\underline{\hspace{#1}}}

% ---------------------------------------------------------------
% Intestazione
% ---------------------------------------------------------------
\title{Pairs trading intraday sull'S\&P~500\\con un agente di notizie}
\author{}   % TODO: nomi del gruppo
\date{\today}

\begin{document}
\maketitle

% ---------------------------------------------------------------
\section{Idea}
% ---------------------------------------------------------------
L'idea alla base del progetto è il \emph{pairs trading}: si considerano le azioni di due aziende con business simili (ad esempio Coca-Cola e Pepsi). Se i due business sono simili, ci si aspetta che anche l'andamento delle loro azioni lo sia. L'idea è quindi di sfruttare i momenti in cui i due titoli si discostano l'uno dall'altro: si acquista (posizione \emph{long}) il titolo rimasto indietro e si vende allo scoperto (posizione \emph{short}) quello che sta correndo di più, scommettendo sul fatto che i prezzi torneranno a riavvicinarsi.

Abbiamo scelto di lavorare su dati intraday a 1 minuto per cogliere anche le variazioni più piccole: con frequenze più basse c'è il rischio di perdere delle opportunità, perché le finestre in cui i prezzi si discostano potrebbero essere molto brevi.

Abbiamo poi deciso di aggiungere un livello ulteriore per rafforzare il segnale: un agente che controlla le notizie al momento dello scostamento di prezzo. Se la notizia è rilevante, è più probabile che i titoli non tornino presto a riavvicinarsi; se invece non c'è nessuna notizia, lo scostamento è probabilmente un errore del mercato e i prezzi si riaggiusteranno.

Il progetto si articola quindi in tre fasi principali:
\begin{itemize}
  \item scelta delle coppie;
  \item calcolo dello spread tra le due azioni;
  \item implementazione dell'agente.
\end{itemize}

% ---------------------------------------------------------------
\section{Raccolta dei dati}
% ---------------------------------------------------------------
I dati di prezzo sono stati scaricati dal terminale Bloomberg, in due forme.

\begin{itemize}
  \item \emph{Dati giornalieri}: cinque anni, dal 16 giugno 2021 al 15 giugno 2026. Sono prezzi di chiusura aggiustati per split e dividendi, accompagnati dai volumi, che servono per il filtro di liquidità. Sono stati usati solo per la scelta delle coppie.
  \item \emph{Dati intraday}: barre da 1 minuto, dal 16 giugno al 30 ottobre 2026, per un totale di 96 sedute di borsa. Sono quelli su cui si opera. Riguardano solo i 30 titoli delle 15 coppie scelte. Quando un titolo non ha scambi in un minuto, la barra viene riempita con il prezzo precedente e segnata come tale: su quei minuti non si apre nessuna operazione.
\end{itemize}

Il periodo intraday è diviso in due parti: il \emph{train}, fino al 16 settembre 2026, su cui abbiamo costruito e confrontato le versioni del modello, e il \emph{test}, dal 17 settembre, che è stato usato una sola volta alla fine.

Le notizie sono state scaricate tramite l'API gratuita per le notizie finanziarie di Alpaca (\url{https://alpaca.markets/}), la cui fonte è Benzinga. Abbiamo scaricato 2.106 notizie per il periodo di train e 2.725 per il periodo di test.

Per il confronto con il mercato (benchmark) abbiamo scaricato anche i prezzi giornalieri di SPY, l'ETF che replica l'S\&P~500, sempre da Alpaca.

% ---------------------------------------------------------------
\section{Scelta delle coppie}
% ---------------------------------------------------------------
La scelta usa soltanto dati giornalieri, su una finestra di cinque anni: dal 16 giugno 2021 al 15 giugno 2026. Si finisce il 15 giugno perché i dati intraday su cui si opera iniziano dopo, così la selezione non vede mai il periodo di trading.

Le coppie sono state scelte secondo i seguenti parametri:
\begin{itemize}
  \item correlazione dei rendimenti giornalieri compresa tra 0,6 e 0,95. Se due titoli sono correlati, i loro rendimenti tendono a salire e scendere insieme: è il legame che ci serve per aspettarci che, dopo uno scostamento, i prezzi tornino a riavvicinarsi. Sotto 0,6 questo legame è troppo debole. Una correlazione troppo alta ($>0{,}95$), invece, indica probabilmente che i due titoli appartengono alla stessa azienda (ad esempio le due classi di azioni di Google, GOOGL e GOOG): non sarebbero una vera coppia di aziende diverse, e lo spread si muoverebbe pochissimo. Una correlazione alta, però, non garantisce che lo spread torni alla media (vedi il punto su half-life e cointegrazione più sotto);
  \item SSD (\emph{Sum of Squared Differences}, somma dei quadrati delle differenze) tra i prezzi normalizzati. I prezzi giornalieri dei due titoli vengono normalizzati, facendoli partire entrambi da 1, in modo da confrontare i percorsi e non i livelli di prezzo. Per ogni giorno si calcola la differenza tra i due percorsi, la si eleva al quadrato (così le differenze positive e negative non si compensano e quelle grandi pesano di più) e si sommano tutti i giorni. Più la SSD è bassa, più i due titoli sono rimasti vicini nel periodo: le coppie vengono ordinate per SSD crescente e si scelgono le più basse. È il metodo classico della letteratura sul pairs trading (Gatev et al., 2006). La SSD misura però la vicinanza dei prezzi, non la tendenza a tornare alla media: due titoli possono essere rimasti vicini per caso (vedi il paragrafo su half-life e cointegrazione più sotto);
  \item stessa sotto-industria GICS, usando i raggruppamenti forniti da Bloomberg. Aziende della stessa sotto-industria hanno lo stesso tipo di business e sono esposte agli stessi fattori economici (domanda, costi, regolamentazione), quindi è ragionevole aspettarsi che i loro prezzi reagiscano in modo simile agli stessi eventi. Questo criterio dà alla scelta una motivazione economica, oltre a quelle statistiche;
  \item liquidità: valore scambiato giornaliero mediano (prezzo per volume) di almeno 100 milioni di dollari. Serve per assicurarci che ci siano abbastanza transazioni su quel titolo, così da poter comprare e vendere senza muovere il prezzo. La soglia è una scelta del gruppo, non deriva da una fonte ufficiale;
  \item dati completi: al massimo l'1\% di prezzi giornalieri mancanti nel periodo di selezione. Con dei buchi la SSD verrebbe calcolata su giorni diversi da titolo a titolo, oppure su prezzi inventati (riempiti con l'ultimo prezzo noto), e la distanza tra le coppie risulterebbe falsata;
  \item un titolo per coppia (lo stesso titolo non compare in più coppie);
  \item al massimo tre coppie per gruppo industriale;
  \item 15 coppie in totale.
\end{itemize}

\paragraph{Procedura.} Si parte dai 501 titoli dell'indice: 166 superano il filtro di liquidità e, tra questi, 154 anche quello dei dati completi. Le coppie possibili della stessa sotto-industria tra questi 154 titoli sono 249. Di queste, 176 vengono scartate perché la correlazione è inferiore a 0,6 e una (GOOG--GOOGL, correlazione 0,996) perché superiore a 0,95: ne restano 72 candidate. Le candidate vengono ordinate per SSD crescente e la classifica si percorre dall'alto, prendendo una coppia solo se nessuno dei suoi titoli è già stato usato e se il suo gruppo industriale non ha già tre coppie, fino ad arrivare a 15. Nessuna coppia entra o esce a mano.

Due misure statistiche non sono state usate per scegliere le coppie, ma solo a posteriori, per verificare se il legame tra i due titoli fosse davvero solido. Entrambe sono calcolate sui dati giornalieri.
\begin{itemize}
  \item \emph{Cointegrazione}: due titoli sono cointegrati se la loro differenza (lo spread) tende a tornare verso un valore medio, invece di allontanarsi senza limite. Questa è esattamente la proprietà su cui si basa il pairs trading. La verifichiamo con il test di Engle-Granger, considerando i due sensi della relazione e tenendo il risultato peggiore. Il legame è considerato statisticamente significativo se il $p$-value è inferiore a 0,05.
  \item \emph{Half-life} (emivita): è il tempo che lo spread impiega, in media, a ricoprire metà della distanza dalla sua media. Più è breve, più il ritorno è rapido e la coppia è utile per operare. Se è infinita, lo spread non mostra alcuna tendenza a tornare alla media.
\end{itemize}
Il risultato è deludente: solo una coppia su 15 supera il test di cointegrazione ($p<0{,}05$), e le half-life vanno da 55 a 157 giorni, con due coppie (KO-PEP e BAC-C) per cui è infinita. Il legame statistico si è quindi rivelato debole. Le coppie vanno di conseguenza valutate sul campo: l'accoppiamento è un'ipotesi, non un fatto verificato.

% ---------------------------------------------------------------
\section{Calcolo dello spread tra le azioni di una coppia}
% ---------------------------------------------------------------
Lo spread tra le due azioni di una coppia viene calcolato a ogni minuto con la formula
\begin{equation}
  S_t = \ln(P^{A}_t) - \ln(P^{B}_t),
\end{equation}
dove $A$ e $B$ sono i due titoli della coppia. Usiamo i logaritmi perché così una variazione di $S$ è, in prima approssimazione, la differenza tra i rendimenti dei due titoli, indipendentemente dal livello dei prezzi. Su ciascuna delle due posizioni della coppia (le due \emph{gambe}, una lunga e una corta) investiamo lo stesso importo (10.000 dollari), quindi il rapporto di copertura tra i due titoli ($\beta$) è uguale a 1.

\paragraph{Z-score.} Per capire se lo spread è nella norma, lo confrontiamo con la sua media recente, misurando di quante deviazioni standard se ne discosta:
\begin{equation}
  Z_t = \frac{S_t - \text{media esponenziale di } S}{\text{deviazione standard esponenziale di } S}.
\end{equation}
Media e deviazione standard sono \emph{esponenziali}: i minuti più recenti pesano di più di quelli lontani, con un tempo di dimezzamento dei pesi di 60 minuti (il peso di un minuto si dimezza dopo 60 minuti). Il calcolo continua anche da un giorno all'altro, senza ripartire ogni mattina, e usa solo dati passati: a ogni minuto si usano solo i prezzi fino a quel minuto. Prima di fidarsi del valore servono almeno 60 minuti di storia.

\paragraph{Segnale di ingresso.} Se $|Z_t| \ge 2$ si considera che ci sia un segnale. Se $Z_t \ge 2$ lo spread è troppo alto: si vende allo scoperto $A$ e si compra $B$. Se $Z_t \le -2$ è troppo basso: si compra $A$ e si vende allo scoperto $B$. Il segnale si legge alla chiusura del minuto $t$ e l'operazione si esegue all'apertura del minuto $t+1$, per non usare informazioni che non sarebbero ancora disponibili. Non si entra nei primi 15 e negli ultimi 30 minuti della seduta, né quando uno dei due titoli ha il prezzo "vecchio" (barra senza scambi, riempita con il prezzo precedente).

\paragraph{Regole di uscita.} Una volta aperta, l'operazione si chiude al primo di questi eventi:
\begin{itemize}
  \item \emph{ritorno alla media}: $Z_t$ torna a 0, cioè lo spread è tornato alla sua media. È l'esito che ci aspettiamo;
  \item \emph{stop}: $|Z_t| \ge 3{,}5$, cioè lo spread si è allontanato di altre 1,5 deviazioni standard oltre la soglia di ingresso. In questo caso si limita la perdita;
  \item \emph{tempo}: sono passati 90 minuti;
  \item \emph{fine giornata}: l'operazione viene chiusa alla chiusura della seduta, senza tenere posizioni overnight.
\end{itemize}
L'uscita avviene all'apertura del minuto successivo a quello in cui la regola scatta, ad eccezione della chiusura di fine giornata.

% ---------------------------------------------------------------
\section{Implementazione dell'agente}
% ---------------------------------------------------------------
L'agente interviene a ogni segnale, cioè ogni volta che $|Z_t| \ge 2$, e legge le notizie uscite prima di quel momento. La sua idea è che uno spread che si allarga per un \emph{fatto strutturale} su una sola delle due aziende (un'acquisizione, un'indagine, una frode, un cambio del management) possa non rientrare, mentre uno che si allarga senza notizie rilevanti sia rumore e rientri. Per ogni segnale l'agente restituisce una decisione, e da qui nascono le tre versioni che confrontiamo:
\begin{itemize}
  \item \textbf{A} (senza agente): si entra a ogni segnale. È il riferimento;
  \item \textbf{B} (l'agente blocca): si entra a ogni segnale, tranne quando l'agente trova un fatto strutturale su una sola delle due aziende; in quel caso si sta fermi;
  \item \textbf{C} (l'agente blocca o segue): come la B, ma in alcuni casi, invece di stare fermi, si entra nel verso della notizia, cioè con la scommessa opposta a quella della A. Il trade esce quando lo spread rientra, dopo 90 minuti o a fine giornata, senza stop.
\end{itemize}

\paragraph{L'agente scelto: FinBERT.} Sul train abbiamo provato diverse varianti dell'agente (vedi sotto) e per il test abbiamo scelto \texttt{finbert\_episodio}, che usa FinBERT. FinBERT è un modello linguistico di dimensioni ridotte, addestrato a riconoscere il sentimento dei testi finanziari: per ogni notizia, dalla lettura del titolo e del sommario, restituisce la probabilità che sia positiva, negativa o neutra. Non risponde a domande e non capisce, ad esempio, chi sia il bersaglio di un'acquisizione: per questo il tipo di evento viene riconosciuto cercando nel testo alcune parole (come ``acquisizione'', ``indagine'', ``dimissioni''). Il modello è inoltre stato pubblicato nel 2020, ben prima del nostro periodo di trading, quindi non può conoscere come sono andati gli eventi. L'agente funziona così:
\begin{itemize}
  \item \emph{Finestra di notizie}: sono le notizie uscite mentre lo spread si allontanava dalla media. La finestra parte dalla prima barra di una serie continua in cui $|Z|$ resta sopra 1 con lo stesso segno (meno 5 minuti) e arriva fino al segnale; se la serie parte dall'apertura, si parte dalle 16:00 del giorno di borsa precedente;
  \item \emph{Notizie rilevanti}: una notizia è \emph{mirata} se cita una sola delle due aziende e al massimo 3 titoli in tutto (si escludono così gli elenchi di titoli), ed è \emph{forte} se la classe che FinBERT giudica più probabile ha probabilità di almeno 0,8;
  \item \emph{Versione B}: si blocca se c'è almeno una notizia mirata e forte, se tutte quelle forti riguardano la stessa azienda e se almeno una contiene una parola che indica un evento. Altrimenti si entra;
  \item \emph{Versione C}: si segue la notizia se la differenza tra i punteggi medi delle due aziende (azienda A meno azienda B) ha modulo almeno 0,5 e lo stesso segno del movimento relativo dei 30 minuti precedenti il segnale; altrimenti si decide come nella B.
\end{itemize}
Le soglie (0,8 e 0,5) e la lista di parole sono state fissate sul train, prima di guardare il test.

\paragraph{Le altre varianti.} Per scegliere l'agente abbiamo provato sul train 21 varianti, confrontate tutte con la stessa A (o con una A ritardata, quando la variante entra in ritardo). Le principali:
\begin{itemize}
  \item \emph{Con un modello linguistico generativo}: Spark X2.6 4B, che gira in locale e legge una notizia alla volta, rispondendo a chi si riferisce, di che tipo è e con quale direzione; le regole di decisione sono scritte nel codice, non sono il modello a deciderle. È stato provato con notizie delle 24 e delle 6 ore precedenti; i nomi delle aziende erano sostituiti con ``Company A'' e ``Company B'' per ridurre il rischio che il modello conoscesse già come era andata a finire. Su un campione di 40 segnali classificati a mano dal gruppo, le decisioni coincidono in 38 casi su 40 (95\%). Una variante analoga usa Llama 3.2 3B;
  \item \emph{Senza modello}: un dizionario di parole positive e negative della finanza (Loughran-\hspace{0pt}McDonald);
  \item \emph{Con FinBERT}: oltre alla variante scelta, una con la finestra legata alla seduta e una che segue solo notizie uscite nei 60 minuti precedenti il segnale;
  \item \emph{Con regole sul testo}: per le notizie di acquisizioni, si blocca solo se l'acquisizione è definitiva e uno dei due titoli è il bersaglio;
  \item \emph{Cambiando il modo di usare la decisione}: ridurre a metà la dimensione dell'operazione invece di bloccarla, oppure entrare 5 o 15 minuti dopo il segnale.
\end{itemize}
Sono state scartate come candidate per il test le varianti la cui regola era stata scritta dopo aver visto i risultati del train e quelle che facevano gli stessi trade della A, per cui la regola non scattava mai. Tra le candidate rimaste abbiamo scelto \texttt{finbert\_episodio} perché entra al momento del segnale come la A, per cui il confronto sul test è diretto.

% ---------------------------------------------------------------
\section{Risultati}
% ---------------------------------------------------------------
I risultati riguardano due periodi: il \emph{train} (dal 16 giugno al 16 settembre 2026, 64 sedute), su cui abbiamo costruito e confrontato le varianti, e il \emph{test} (dal 17 settembre al 30 ottobre, 32 sedute), usato una sola volta per la variante scelta. Il capitale è di 300.000 dollari, 10.000 per ciascuna delle due posizioni di ogni coppia. I guadagni sono espressi in punti base (bps): 1~bps è lo 0,01\% dell'importo investito, cioè 1 dollaro ogni 10.000. Il costo di transazione è di 0 bps nello scenario principale, come da indicazione del docente, e di 5 bps per operazione in quello realistico; ogni trade ha 4 operazioni (aprire e chiudere le due posizioni).

\paragraph{La strategia di base.} Senza agente (versione A) lo spread rientra e la strategia guadagna, ma poco. Nel train fa 973 operazioni con un guadagno medio di 7,4~bps ciascuna, cioè 7.239 dollari (+2,4\% del capitale); nel test 500 operazioni, 4,9~bps ciascuna, 2.472 dollari (+0,8\%). Il guadagno non sopravvive ai costi: il costo per operazione a cui il guadagno medio si azzera è di 1,9~bps nel train e di 1,2~bps nel test. Con 5~bps la strategia perde 12.221 dollari nel train (--4,1\%) e 7.528 dollari nel test (--2,5\%).

\paragraph{Gli agenti.} La tabella riporta il test, a costo~0, per le tre versioni della variante scelta, \texttt{finbert\_episodio}.
\begin{table}[h]
\centering\small
\begin{tabular}{lrrrr}
\toprule
 & Operazioni & Medio (bps) & Totale (\$) & Diff. dalla A (\$) \\
\midrule
A (senza agente)      & 500 & 4,9 & 2.472 & \\
B (l'agente blocca)   & 493 & 4,6 & 2.284 & --189 \\
C (blocca o segue)    & 495 & 2,5 & 1.230 & --1.242 \\
\bottomrule
\end{tabular}
\end{table}
L'agente non migliora la strategia. Nella versione B la differenza dalla A è negativa e piccola: ricampionando i giorni con reinserimento (10.000 volte), l'intervallo di confidenza al 95\% va da --527 a +135 dollari, quindi comprende lo zero e la B non si distingue dalla A. L'agente blocca soltanto 7 operazioni, che in media avrebbero guadagnato 27~bps contro i 4,6 delle altre, e rientrano quasi come le altre (43\% contro 49\%); un blocco fatto a caso dello stesso numero di operazioni avrebbe fatto meglio nel 90\% dei casi. Nella versione C l'intervallo di confidenza va da --2.183 a --374 dollari, tutto sotto lo zero: seguire la notizia peggiora la strategia. L'agente segue la notizia 21 volte, e la notizia ha ragione (lo spread non rientra) soltanto 5 volte, il 24\%; nel train era il 55\%. Nel train la differenza dalla A aveva lo stesso segno: --285 dollari per la B e --651 per la C. Con un costo di 5~bps le differenze dalla A cambiano poco (--49 dollari per la B e --1.142 per la C): i costi si aggiungono a ogni operazione e non dicono nulla sulla qualità delle decisioni dell'agente, quindi non cambiano il giudizio.

\paragraph{Le altre varianti.} Sul train nessuna delle 21 varianti ha battuto la A: le B perdono tra circa 60 e 1.530 dollari rispetto alla A, le C fanno peggio delle B con le stesse decisioni. Anche il dizionario di parole positive e negative, senza alcun modello, blocca come farebbe il caso.

\paragraph{Il confronto con il mercato.} La tabella confronta il rendimento cumulato delle tre versioni con quello di SPY (l'S\&P~500 comprato e tenuto).
\begin{table}[h]
\centering\small
\begin{tabular}{lrrrr}
\toprule
 & A & B & C & SPY \\
\midrule
Train, costo 0      & +2,42\% & +2,32\% & +2,20\% & +0,15\% \\
Train, costo 5~bps  & --4,07\% & --4,13\% & --4,26\% & \\
Test, costo 0       & +0,83\% & +0,76\% & +0,41\% & +0,47\% \\
Test, costo 5~bps   & --2,51\% & --2,52\% & --2,89\% & \\
\bottomrule
\end{tabular}
\end{table}
Senza costi, la strategia fa meglio di SPY nel train in tutte e tre le versioni; nel test fanno meglio la A e la B, mentre la C fa leggermente peggio. Con i costi, invece, tutte e tre perdono. La strategia oscilla molto meno della borsa (volatilità annua intorno all'1\% contro l'11\% di SPY) perché è investita solo di giorno e in poche operazioni alla volta, per cui gli indicatori come lo Sharpe ratio non sono confrontabili con quelli di un fondo e, su soli 32 giorni, sono poco affidabili. La strategia è inoltre neutrale rispetto al mercato: nel test il beta rispetto a SPY è --0,04 e, nei 16 giorni in cui SPY scende, la A guadagna in media 4,0~bps al giorno.

\paragraph{In sintesi.} Lo spread tra le coppie rientra e la strategia di base ha un piccolo guadagno lordo, che con costi realistici diventa una perdita. L'agente di notizie non migliora la strategia: nel test, dove è stato provato una sola volta, bloccare le operazioni non si distingue dal non farlo, e seguire la notizia peggiora il risultato. Questi risultati poggiano su pochi casi (7 blocchi e 21 notizie seguite nel test) e vanno letti con prudenza.

% ---------------------------------------------------------------
\section{Limiti}
% ---------------------------------------------------------------
I limiti principali sono:
\begin{itemize}
  \item la scarsità di dati: Bloomberg fornisce le barre da 1 minuto solo per gli ultimi mesi, quindi abbiamo 96 sedute (64 di train e 32 di test), che sono tutta la storia intraday disponibile. In tutto il periodo l'S\&P~500 è quasi fermo (+0,15\% nel train, +0,47\% nel test) e con volatilità simile: i risultati valgono per questo tipo di mercato e non sappiamo come andrebbero in un mercato in forte calo o molto volatile;
  \item il numero ridotto di casi per l'agente: nel test blocca 7 operazioni e segue 21 notizie, quindi ogni conclusione sull'agente è fragile;
  \item costi stimati e non misurati, vendita allo scoperto sempre possibile e senza costi di prestito, ordini eseguiti senza effetto sul prezzo;
  \item prezzi intraday non aggiustati per i dividendi, il cui stacco può spostare lo spread all'apertura;
  \item una sola fonte di notizie, con una copertura che può essere incompleta.
\end{itemize}

\paragraph{Errori tipici del backtest.} Il corso chiede di mostrare come il metodo affronta cinque errori tipici.
\begin{itemize}
  \item \emph{Overfitting}: i parametri del segnale sono stati scelti una volta, con una griglia di 32 combinazioni sul train delle prime 10 coppie e una regola di scelta scritta prima, e poi tenuti uguali per tutte le 15 coppie. Le soglie dell'agente (0,8 per una notizia forte, 0,5 per seguire la notizia, al massimo 3 titoli) sono le stesse sul train e sul test. Il test è stato usato una sola volta, con le regole bloccate;
  \item \emph{Lookahead}: il segnale si calcola con i prezzi fino alla chiusura del minuto e si esegue all'apertura del minuto successivo. Le notizie si filtrano con l'orario di pubblicazione, e l'agente vede solo quelle uscite prima della decisione. Le coppie sono scelte con dati giornalieri fino al 15 giugno, prima dell'inizio dei dati intraday. Resta il fatto che la classificazione GICS fornita da Bloomberg è quella attuale e non quella storica;
  \item \emph{Survivorship}: l'universo sono i titoli dell'S\&P~500 di oggi, quindi chi è uscito dall'indice manca. Le coppie sono scelte tra aziende che sappiamo essere rimaste nell'indice;
  \item \emph{Data snooping}: sul train abbiamo provato 21 varianti dell'agente, quindi la migliore lo è in parte anche per caso. Il criterio con cui abbiamo scelto la variante per il test è stato precisato dopo aver visto il train, ma prima del test; le regole del test sono state scritte prima del giro ma messe su git dopo. Anche le soglie della selezione delle coppie (liquidità, correlazione, tetto per gruppo, numero di coppie) sono scelte del gruppo, e alcune sono state fissate dopo una prima esplorazione dei dati. Per distinguere la fortuna dalla differenza reale abbiamo usato il ricampionamento dei giorni, il confronto con blocchi fatti a caso e il test di Fisher;
  \item \emph{Leakage dei modelli preaddestrati}: la variante portata sul test usa FinBERT, pubblicato tra il 2019 e il 2020 e quindi senza conoscenza degli eventi del 2026; vede solo titolo e sommario delle notizie. I modelli generativi delle altre varianti (Spark X2.6 4B e Llama 3.2 3B) sono stati usati solo sul train, con i nomi delle aziende nascosti.
\end{itemize}

% ---------------------------------------------------------------
\section{Proposte}
% ---------------------------------------------------------------
Dai limiti e dai risultati nascono alcune proposte di miglioramento.
\begin{itemize}
  \item \emph{Più dati e regimi di mercato diversi.} Con un periodo più lungo, o con dati intraday di più anni, si potrebbe valutare la strategia in mercati in calo o molto volatili, e l'agente avrebbe più casi su cui essere giudicato (nel test blocca soltanto 7 operazioni);
  \item \emph{Isolare l'effetto su un singolo titolo.} Attualmente si ipotizza che lo scostamento derivi dalla variazione di entrambe le azioni e non di una sola. Non si isola quindi l'effetto su un singolo titolo: quando c'è un segnale si compra un titolo e si vende l'altro, anche se magari a essere «fuori trend» è soltanto uno dei due;
  \item \emph{Una scelta delle coppie più solida.} Il legame statistico tra i titoli si è rivelato debole: una sola coppia su 15 supera il test di cointegrazione. Si potrebbero usare la cointegrazione e l'half-life come criteri di scelta, e non solo come verifica a posteriori;
  \item \emph{Un agente più potente.} Con un modello accessibile tramite API e con più fonti di notizie, l'agente potrebbe riconoscere meglio i fatti che davvero impediscono allo spread di rientrare, distinguendoli dai comunicati ripetitivi (ad esempio gli studi legali che annunciano indagini) che nel test hanno prodotto molti blocchi;
  \item \emph{Una validazione separata per scegliere la variante.} Oggi la variante è scelta su tutto il train; con una parte dedicata alla sola scelta si ridurrebbe il rischio che la migliore lo sia per caso;
  \item \emph{Costi misurati.} Una stima dei costi reali di negoziazione e di prestito dei titoli venduti allo scoperto permetterebbe di capire se esista un livello di costo a cui la strategia resta in guadagno.
\end{itemize}

\end{document}
